\section{Prime supplies and the quantitative target}\label{q:setup}

Fix integers \(h,l\ge1\) and \(b\in\Z\). Throughout Part~I, constants
in estimates may depend on these fixed data unless declared absolute.
We will prove the following intermediate statement.

\begin{proposition}[Quantitative correlations in every progression]
\label{q:progression}
There is one absolute \(c>0\) such that, for every fixed \(h,l\ge1\)
and \(b\in\Z\),
\[
 F(X):=\frac1X\sum_{1\le n\le X}
   \1_{n\equiv b\pmod l}\lambda(n)\lambda(n+h)
 \ll_{h,l}(\log X)^{-c}\qquad(X\ge3).
\]
The estimate holds for every real \(X\), with no condition on
\(\gcd(b,l)\).
\end{proposition}

There are only \(l\) residue classes, so the implied constant can be
chosen independently of the representative \(b\).
It suffices to work above a threshold depending on the fixed data:
bounded \(X\) can be absorbed into that constant.
Our first task is to find many divisors in short multiplicative intervals,
retaining a fixed positive fraction of their total reciprocal weight.
The construction separates primes used for centering from primes used
to supply those divisors.

The proof uses two absolute constants. First choose \(A\ge1000\)
large enough for the finite-law comparison in
\Cref{q:finite-law}. Later choose \(W\ge10\) after the absolute
graph constants have been established. Set
\begin{equation}\label{q:parameters}
 \begin{gathered}
 L=(\log X)^{1/A},\qquad \delta=\frac1{200},\qquad
 J=\left\lfloor\frac{\delta\log L}{6W}\right\rfloor,\\
 Y_i=L^{1-\delta}e^{6Wi}\quad(0\le i\le J),\qquad
 H_0=e^{L^{1-\delta}},\qquad
 \eta=e^{-J},\qquad K=e^{4J}.
 \end{gathered}
\end{equation}
We may assume \(J\ge1\). Since \(6WJ\le\delta\log L\), one has
\(Y_J\le L\).

\subsection{Centered primes and padding primes}

For \(1\le i\le J\), let
\[
 P_i=\{p:p\equiv1\pmod5,\quad Y_{i-1}\le\log p<Y_i,\quad
                              p\nmid lh\}.
\]
Let \(P=\bigcup_iP_i\), and let
\[
 Q=\{p:p\not\equiv1\pmod5,\quad p\le e^L,\quad p\nmid lh\}.
\]
These sets are disjoint, and every prime in \(P\) is at least \(H_0\).
Write
\[
 V_i=\sum_{p\in P_i}\frac1p,\qquad V_*=\prod_{i=1}^J V_i,\qquad
 \mathcal D=\left\{\prod_{i=1}^J p_i:p_i\in P_i\right\}.
\]
A label \(d\in\mathcal D\) determines its ordered tuple of primes,
because the supplies are disjoint.

The prime number theorem for the fixed modulus \(5\)
\cite[Theorem~12.1]{Koukoulopoulos19}, followed by removal of the
finitely many primes dividing \(lh\), gives
\begin{equation}\label{q:pnt}
 \sum_{\substack{p\le x\\p\in\mathcal R}}\log p
 =\alpha_{\mathcal R}x+
 O_{h,l}\!\left(xe^{-c_1\sqrt{\log x}}\right),
\end{equation}
where \(\mathcal R\) is either the primes congruent to \(1\pmod5\)
or the primes not congruent to \(1\pmod5\), in either case excluding
prime divisors of \(lh\), and
\(\alpha_{\mathcal R}=1/4\) or \(3/4\), respectively.
Here \(c_1>0\) is absolute; the threshold and the implied constant
may depend on \(h,l\). In the second selection the possible prime
\(5\) contributes only a bounded term.

Partial summation on the logarithmic interval
\([Y_{i-1},Y_i)\) gives
\[
 V_i=\frac14\log\frac{Y_i}{Y_{i-1}}+o(1)
     =\frac32W+o(1),
\]
uniformly for \(1\le i\le J\). Indeed all these intervals begin at
least at \(L^{1-\delta}\), so their accumulated endpoint and
integrated error in \eqref{q:pnt} tends to zero uniformly.
Consequently, for sufficiently large \(X\),
\begin{equation}\label{q:prime-masses}
 W\le V_i\le2W,\qquad
 \log d\le\sum_{i=1}^JY_i
 \le \frac{L}{1-e^{-6W}}<2L\quad(d\in\mathcal D).
\end{equation}

Let \(\mathcal Q\) be the squarefree products of primes in \(Q\),
including \(1\), and put
\begin{equation}\label{q:padding-weights}
 u(q)=4^{\omega(q)},\qquad
 S=\sum_{q\in\mathcal Q}\frac{u(q)}q
   =\prod_{p\in Q}(1+4/p),\qquad
 w(n)=\sum_{\substack{q\in\mathcal Q\\q\mid n}}u(q)
   =5^{\omega_Q(n)}.
\end{equation}
Here \(\omega\) counts distinct prime factors. For any finite prime
set \(\mathcal R\), write
\(\omega_{\mathcal R}(n)=\#\{p\in\mathcal R:p\mid n\}\). The factor \(q\) will be
called the padding factor. All sums over these finite prime sets
are finite, however large their cardinalities.

\subsection{Short multiplicative intervals}

For every integer \(j\ge0\) with \(j\eta\le100L\), define
\begin{equation}\label{q:bins}
 \mathcal S_j=
 \{(d,q)\in\mathcal D\times\mathcal Q:
    \omega(q)\le100\log L,\quad
    j\eta\le\log(dq)<(j+1)\eta\},
 \qquad T_j=Xe^{j\eta}.
\end{equation}
There are \(O(L/\eta)\) such intervals. Define their total harmonic
mass by
\[
 S_0=\sum_j\sum_{(d,q)\in\mathcal S_j}\frac{u(q)}{dq}.
\]
\begin{lemma}[Mass retained in the intervals]\label{q:bin-mass}
For sufficiently large \(X\),
\[
 \frac12SV_*\le S_0\le SV_*.
\]
\end{lemma}
\begin{proof}
Under the law \(u(q)/(qS)\), primes in \(Q\) are independently
included with probabilities \(4/(p+4)\).
Equation \eqref{q:pnt} and partial summation give
\[
 \mathbb E_q\log q\le4L,\qquad
 \mathbb E_q\omega(q)\le4\log L
\]
for large \(L\). Markov's inequality therefore excludes
\(\log q>98L\) with probability at most \(4/98\), and excludes
\(\omega(q)>100\log L\) with probability at most \(4/100\).
Every remaining \(q\), together with every \(d\in\mathcal D\),
satisfies \(\log(dq)\le100L\) by \eqref{q:prime-masses}, and so lies
in one of the intervals \eqref{q:bins}. Their probability exceeds
\(1/2\). Summing the independent harmonic weights of \(d\), whose
total is \(V_*\), proves the lower bound. Dropping both restrictions
proves the upper bound.
\end{proof}

\subsection{Centered divisor sums}

The bins let us compare dilated correlations at essentially the same
cutoff. For each bin define
\begin{equation}\label{q:centered-average}
\begin{split}
 C_j={}&\frac1{T_j}\sum_{1\le n\le T_j}
 \sum_{(d,q)\in\mathcal S_j}
 u(q)\1_{q\mid n}\1_{n\equiv bqd\pmod l}\\
 &\hspace{13mm}\cdot
 \prod_{p\mid d}\left(\1_{p\mid n}-\frac1p\right)
 \lambda(n)\lambda(n+hqd).
\end{split}
\end{equation}
For a fixed pair $(d,q)$, the term retaining every divisibility indicator
in this expansion forces $qd\mid n$. The substitution $n=qdm$ turns it into
\[
 \frac{u(q)}{qd}F\left(\frac{T_j}{qd}\right),
 \qquad Xe^{-\eta}<\frac{T_j}{qd}\le X.
\]
Here $qd$ is coprime to $l$, and complete multiplicativity cancels
$\lambda(qd)^2=1$. Since the summands defining $F$ are bounded by one,
$F(T_j/(qd))=F(X)+O(\eta+X^{-1})$. Summing the weights and using
$S_0\le SV_*$, the full-divisibility contribution is therefore
\begin{equation}\label{q:full-divisibility}
 S_0F(X)+O(SV_*(\eta+X^{-1})).
\end{equation}
\Cref{q:center} will bound the remaining terms. We will then estimate
the centered sums themselves by the graph argument, after removing
summands of small total absolute weight.

\subsection{The independent residue space}

The factor $\1_{p\mid n}-1/p$ has mean zero for a uniform residue
modulo $p$. To use this cancellation simultaneously at many primes,
we introduce the auxiliary product law with one independent uniform
residue modulo each $p\in P\cup Q$, and one independent uniform residue
modulo $l$. The latter is a single coordinate, even if $l$ is not
squarefree; for $l=1$ it is deterministic. The selected primes are all
coprime to $l$. Write $\E$ and $\Pbb$ for expectation and probability
under this law. A translated site $n+a$ uses these same coordinates
shifted by $a$, so divisibility, progression tests, and $w(n+a)$ are
defined without a common integer $n$.

At each site $\E w=S$. The prime sets grow with $X$, so this model
must be compared with finite integer averages before it can be used to
estimate deletion costs and graph moments. The next section supplies
that comparison for the residue tests and truncated weights used below.
